% Generated draft: requires mathematical and source review.
\documentclass[11pt]{article}
\usepackage[margin=25mm]{geometry}
\usepackage{fontspec}
\setmainfont{Latin Modern Roman}
\usepackage{amsmath,amssymb,mathtools,unicode-math}
\setmathfont{Latin Modern Math}
\usepackage{enumitem,needspace,xurl,hyperref}
\setlength{\parindent}{0pt}
\setlength{\parskip}{5pt}
\setlength{\emergencystretch}{3em}
\newcommand{\N}{\mathbb{N}}
\newcommand{\Z}{\mathbb{Z}}
\newcommand{\Q}{\mathbb{Q}}
\newcommand{\R}{\mathbb{R}}
\newcommand{\C}{\mathbb{C}}
\newcommand{\F}{\mathbb{F}}
\newcommand{\A}{\mathbb{A}}
\newcommand{\PP}{\mathbb{P}}
\newcommand{\E}{\mathbb{E}}
\newcommand{\eps}{\varepsilon}
\newcommand{\dd}{\,\mathrm{d}}
\begin{document}
\section*{UnsolvedMath ID: 30006481}
\textbf{Source problem code:} catalogue entry
\subsection{Definitions and mathematical statement}
The supplied record is an incomplete derivation fragment, marked \texttt{statement\_status: unrecoverable}; it does not specify a recoverable standalone problem.

The model description concerns spin-\(\tfrac12\) fermions with spin populations \(N_{\uparrow}\) and \(N_{\downarrow}\), total particle number
\[
N=N_{\uparrow}+N_{\downarrow},
\]
and spatial box
\[
\Lambda=[-L/2,L/2]^3\subset\R^3
\]
with periodic boundary conditions. Here the populations are particle counts and \(L>0\) is the box side length. The interaction \(V_{\infty}\), defined on \(\R^3\), is described as compactly supported and radial. Its stated Lebesgue-space hypothesis is corrupted in the extraction, so no exponent or additional potential hypotheses can be specified reliably.

The fragment introduces a positive scattering length \(a>0\) through the displayed relation
\[
8\pi a=\int_{\R^3}V_{\infty}(x)\bigl(1-\varphi_{\infty}(x)\bigr)\,\dd x.
\]
This is a transcription of the supplied relation, not a complete definition: the record does not define \(\varphi_{\infty}\), give its equation or normalization, or supply sufficient hypotheses to justify the integral and the asserted positivity.

The opening expression for \(F(1)\) has corrupted fractions or exponents, and \(F\) is undefined. The fragment says that this expression agrees with the Huang--Yang conjecture, labeled \((1)\), and announces a discussion of ideas used in reference \([2]\) to prove \((1)\). Neither equation \((1)\) nor an identification of reference \([2]\) is supplied. Consequently, the Hamiltonian, energy quantity, limiting regime, error estimate, and quantified target assertion cannot be recovered from this record.

The supplied status label and dated literature assessment call the record open, whereas the supplied statement-verification note describes it as exposition from a proved Huang--Yang energy expansion rather than an open problem. These are conflicting supplied characterizations; neither constitutes independent verification here. No complete conjecture or theorem is asserted in this draft.
\subsection{Short English statement}
This fragment describes a periodic box containing spin-\(\tfrac12\) fermions and records a scattering-length relation. Missing definitions, corrupted notation, and absent references prevent recovery of a precise Huang--Yang statement. The supplied open label conflicts with a supplied verification note describing the passage as part of a proved expansion.
\subsection{Sources}
\begin{itemize}
\item Local UnsolvedMath registry, permanent ID \(30006481\), title: Huang--Yang Derivation Fragment; exact source problem code: \texttt{catalogue entry}. No proposer is identified. The supplied proposed year is 2025.
\item Supplied source citation: \emph{Arbeitsgemeinschaft: Analysis of Many-body Quantum Systems (2026)}, Oberwolfach Reports, pp.~2491--2544, DOI \texttt{10.4171/owr/2025/46}. \url{https://doi.org/10.4171/owr/2025/46}. The citation's parenthetical year 2026 differs from the report designation 46/2025; the supplied record does not resolve this discrepancy.
\item Supplied literature source: Oberwolfach Report 46/2025. \url{https://doi.org/10.4171/owr/2025/46}.
\item Supplied literature source: \emph{Rigorous upper bound for the Gell--Mann--Brueckner formula}. \url{https://doi.org/10.1007/s00220-023-04672-2}. The record does not explain its relationship to this Huang--Yang fragment or identify it as reference \([2]\).
\item Literature assessment checked at \texttt{2026-08-22}: supplied dated curation, not independently verified. It labels the record open and describes it as a fragment of the Huang--Yang derivation for which no independent closure of the underlying claim was established by that assessment. The background explicitly requests independent verification of time-sensitive status and cited claims.
\item Supplied statement-verification note: the official report is said to identify this extraction as formula and exposition from the derivation of a proved Huang--Yang energy expansion, rather than an open problem. The note says that no separate clean statement was promoted and that the original record was retained publicly for provenance. This verification claim is supplied metadata, not an independent source visit or mathematical assessment here. The record's update timestamp is \texttt{2026-09-16T12:06:53Z}; no separate verification date is given.
\end{itemize}
\end{document}
